% Generated draft: requires mathematical and source review.
\documentclass[11pt]{article}
\usepackage[margin=25mm]{geometry}
\usepackage{fontspec}
\setmainfont{Latin Modern Roman}
\usepackage{amsmath,amssymb,mathtools,unicode-math}
\setmathfont{Latin Modern Math}
\usepackage{enumitem,needspace,xurl,hyperref}
\setlength{\parindent}{0pt}
\setlength{\parskip}{5pt}
\setlength{\emergencystretch}{3em}
\newcommand{\N}{\mathbb{N}}
\newcommand{\Z}{\mathbb{Z}}
\newcommand{\Q}{\mathbb{Q}}
\newcommand{\R}{\mathbb{R}}
\newcommand{\C}{\mathbb{C}}
\newcommand{\F}{\mathbb{F}}
\newcommand{\A}{\mathbb{A}}
\newcommand{\PP}{\mathbb{P}}
\newcommand{\E}{\mathbb{E}}
\newcommand{\eps}{\varepsilon}
\newcommand{\dd}{\,\mathrm{d}}
\begin{document}
\section*{UnsolvedMath ID: 700411}
\textbf{Source problem code:} catalogue entry
\subsection{Definitions and mathematical statement}
Let \(G\) be a group. For subgroups \(H,K\leq G\), write \(\langle H\cup K\rangle\) for the smallest subgroup of \(G\) containing both. Following the source, call \(G\) \emph{indecomposable} if
\[
\forall H,K\leq G,\qquad (H\neq G\ \text{and}\ K\neq G)\ \Longrightarrow\ \langle H\cup K\rangle\neq G.
\]
This is the source-specific meaning of indecomposable. A group is \emph{periodic} if
\[
\forall g\in G\ \exists n\in\Z_{\geq 1}\qquad g^n=1.
\]
Let
\[
G'=\langle x^{-1}y^{-1}xy:x,y\in G\rangle
\]
be the commutator subgroup. The group \(G\) is \emph{metabelian} if \(G'\) is abelian, equivalently \((G')'=\{1\}\).

\textbf{Classification problem.} Describe all periodic metabelian groups that are indecomposable in the above sense. No finiteness, finite generation, or uniform bound on element orders is assumed. The supplied statement does not exclude the trivial group. It requests a description of this class without prescribing a particular form of classification.
\subsection{Short English statement}
Describe all groups in which every element has finite order, the commutator subgroup is abelian, and any two proper subgroups together generate a proper subgroup. Indecomposability has exactly the meaning specified in the source.
\subsection{Sources}
\begin{itemize}
\item Supplied local UnsolvedMath registry record, permanent ID \(700411\), exact source code \texttt{catalogue entry}, titled \emph{Kourovka Notebook Problem 7.5}; upstream revision \texttt{d371ee917ac9647fbb06cd6fc529deec6f3dda82}.
\item Attributed in the supplied record to V. V. Belyaev, proposed in 1980; original problem number 7.5 in \emph{The Kourovka Notebook (Unsolved Problems in Group Theory)}, 7th issue (1980). The supplied background also identifies the 21st issue edition (2026), \texttt{arXiv:1401.0300v48}.
\item The registry places this problem in \emph{Kourovka Notebook (Issues 1--21)}, whose supplied description names E. I. Khukhro and V. D. Mazurov as editors. The registry labels the problem \emph{open}; this is supplied source metadata, not an independently verified assessment.
\end{itemize}
\end{document}
